\section{Cholesterol Metabolism}

\subsection{Overview of Cholesterol}

\begin{definition}
\textbf{Cholesterol} is a 27-carbon sterol that serves as:
\begin{itemize}
    \item Structural component of cell membranes (modulates fluidity)
    \item Precursor for steroid hormones, bile acids, and vitamin D
    \item Component of lipoproteins
\end{itemize}
\end{definition}

\textbf{Sources of cholesterol:}
\begin{itemize}
    \item \textbf{Endogenous synthesis}: $\sim$800-1000 mg/day (primarily liver)
    \item \textbf{Dietary intake}: $\sim$300-500 mg/day
\end{itemize}

\textbf{Fates of cholesterol:}
\begin{itemize}
    \item Membrane synthesis
    \item Bile acid synthesis (primary excretion route)
    \item Steroid hormone synthesis
    \item Vitamin D synthesis
    \item Esterification for storage/transport
\end{itemize}

\subsection{Cholesterol Biosynthesis}

\textbf{Location:} Cytosol and endoplasmic reticulum

\textbf{Primary tissue:} Liver (major site); also intestine, adrenal cortex, gonads

\textbf{Precursor:} Acetyl-CoA (from carbohydrates, amino acids, or fatty acids)

\textbf{The pathway can be divided into four stages:}

\subsubsection{Stage 1: Synthesis of Mevalonate}

\textbf{Step 1: Formation of Acetoacetyl-CoA}
\[
2 \text{ Acetyl-CoA} \xrightarrow{\text{Thiolase}} \text{Acetoacetyl-CoA} + \text{CoA-SH}
\]

\textbf{Step 2: Formation of HMG-CoA}
\[
\text{Acetoacetyl-CoA} + \text{Acetyl-CoA} \xrightarrow{\text{HMG-CoA synthase}} \text{HMG-CoA}
\]
\begin{itemize}
    \item This is the \textbf{cytosolic} HMG-CoA synthase
    \item Different from mitochondrial enzyme (ketogenesis)
\end{itemize}

\textbf{Step 3: Formation of Mevalonate (RATE-LIMITING STEP)}
\[
\text{HMG-CoA} + 2\text{NADPH} + 2\text{H}^+ \xrightarrow{\textbf{HMG-CoA reductase}} \text{Mevalonate} + 2\text{NADP}^+ + \text{CoA-SH}
\]

\begin{tcolorbox}[colback=yellow!10, colframe=yellow!70!black, title=\textbf{HMG-CoA Reductase: Key Regulatory Enzyme}]
\begin{itemize}
    \item \textbf{Rate-limiting enzyme} of cholesterol synthesis
    \item Located in ER membrane
    \item Target of \textbf{statin drugs} (competitive inhibitors)
    \item Regulated by multiple mechanisms (see regulation section)
\end{itemize}
\end{tcolorbox}

\subsubsection{Stage 2: Conversion of Mevalonate to Isopentenyl Pyrophosphate}

Three ATP-dependent reactions convert mevalonate (C6) to isopentenyl pyrophosphate (C5), the basic 5-carbon isoprenoid unit.

\[
\text{Mevalonate} \xrightarrow{3 \text{ ATP}} \text{Isopentenyl pyrophosphate (IPP)} + \text{CO}_2
\]

IPP can isomerize to dimethylallyl pyrophosphate (DMAPP).

\subsubsection{Stage 3: Condensation of Isoprene Units}

\textbf{Formation of Geranyl Pyrophosphate (C10):}
\[
\text{IPP (C5)} + \text{DMAPP (C5)} \rightarrow \text{Geranyl-PP (C10)}
\]

\textbf{Formation of Farnesyl Pyrophosphate (C15):}
\[
\text{Geranyl-PP (C10)} + \text{IPP (C5)} \rightarrow \text{Farnesyl-PP (C15)}
\]

\textbf{Formation of Squalene (C30):}
\[
2 \text{ Farnesyl-PP (C15)} \xrightarrow{\text{Squalene synthase}} \text{Squalene (C30)} + 2\text{PP}_i
\]
\begin{itemize}
    \item Head-to-head condensation
    \item Requires NADPH
    \item First committed step to cholesterol (farnesyl-PP has other fates)
\end{itemize}

\subsubsection{Stage 4: Cyclization of Squalene to Cholesterol}

\textbf{Squalene Epoxidation:}
\[
\text{Squalene} + \text{O}_2 + \text{NADPH} \xrightarrow{\text{Squalene monooxygenase}} \text{Squalene-2,3-epoxide}
\]

\textbf{Cyclization:}
\[
\text{Squalene-2,3-epoxide} \xrightarrow{\text{Oxidos
qualene cyclase}} \text{Lanosterol}
\]

\textbf{Conversion to Cholesterol:}
\[
\text{Lanosterol} \xrightarrow{\text{19 reactions}} \text{Cholesterol}
\]
\begin{itemize}
    \item Removal of 3 methyl groups
    \item Reduction of double bond in side chain
    \item Migration of double bond from C8 to C5
\end{itemize}

\subsection{Regulation of Cholesterol Synthesis}

Cholesterol synthesis is tightly regulated at multiple levels, primarily targeting HMG-CoA reductase.

\textbf{1. Transcriptional Regulation (SREBP System):}

\begin{tcolorbox}[colback=green!5, colframe=green!70!black, title=\textbf{SREBP-2 Pathway}]
\textbf{When cholesterol is LOW:}
\begin{enumerate}
    \item SREBP-2 (Sterol Regulatory Element Binding Protein) is in ER membrane
    \item SCAP (SREBP Cleavage-Activating Protein) escorts SREBP to Golgi
    \item Site-1 and Site-2 proteases cleave SREBP
    \item Active SREBP fragment enters nucleus
    \item Activates transcription of HMG-CoA reductase and LDL receptor genes
\end{enumerate}

\textbf{When cholesterol is HIGH:}
\begin{enumerate}
    \item Cholesterol binds to SCAP
    \item SCAP conformation changes; binds to Insig (ER retention protein)
    \item SREBP-SCAP complex retained in ER
    \item No transcription of cholesterol synthesis genes
\end{enumerate}
\end{tcolorbox}

\textbf{2. Translational Regulation:}
\begin{itemize}
    \item High cholesterol decreases HMG-CoA reductase mRNA translation
    \item Mediated by sequences in the 5' UTR of mRNA
\end{itemize}

\textbf{3. Protein Degradation:}
\begin{itemize}
    \item High cholesterol/oxysterols accelerate degradation of HMG-CoA reductase
    \item Insig-mediated ubiquitination and proteasomal degradation
    \item Half-life can decrease from >12 hours to <1 hour
\end{itemize}

\textbf{4. Covalent Modification:}
\begin{itemize}
    \item \textbf{AMPK} phosphorylates HMG-CoA reductase $\rightarrow$ \textbf{inactive}
    \item Occurs when cellular energy is low
    \item Prevents ATP consumption for cholesterol synthesis during energy deficit
\end{itemize}

\textbf{5. Hormonal Regulation:}
\begin{itemize}
    \item \textbf{Insulin}: Activates phosphatase $\rightarrow$ dephosphorylates enzyme $\rightarrow$ active
    \item \textbf{Glucagon}: Activates AMPK $\rightarrow$ phosphorylates enzyme $\rightarrow$ inactive
    \item \textbf{Thyroid hormone}: Increases LDL receptor expression
\end{itemize}

\subsection{Cholesterol Transport: Lipoproteins}

\begin{definition}
\textbf{Lipoproteins} are spherical particles that transport hydrophobic lipids (cholesterol, triglycerides) through the aqueous bloodstream. They consist of a hydrophobic core surrounded by an amphipathic shell.
\end{definition}

\textbf{Structure of Lipoproteins:}
\begin{itemize}
    \item \textbf{Core}: Triglycerides and cholesteryl esters (hydrophobic)
    \item \textbf{Shell}: Phospholipids, free cholesterol, apolipoproteins (amphipathic)
\end{itemize}

\begin{center}
\begin{tabular}{lccccc}
\toprule
\textbf{Lipoprotein} & \textbf{Density} & \textbf{Main Lipid} & \textbf{Source} & \textbf{Key Apo} & \textbf{Function} \\
\midrule
Chylomicron & Lowest & TG (dietary) & Intestine & B-48, C-II, E & Dietary lipid transport \\
VLDL & Very low & TG (hepatic) & Liver & B-100, C-II, E & Hepatic TG export \\
IDL & Intermediate & TG/Chol & VLDL remnant & B-100, E & VLDL $\rightarrow$ LDL \\
LDL & Low & Cholesterol & IDL & B-100 & Chol delivery to tissues \\
HDL & Highest & Cholesterol & Liver/Intestine & A-I, A-II & Reverse chol transport \\
\bottomrule
\end{tabular}
\end{center}

\subsection{Lipoprotein Metabolism Pathways}

\subsubsection{Exogenous Pathway (Dietary Lipids)}

\begin{enumerate}
    \item \textbf{Chylomicron assembly} in intestinal enterocytes
    \begin{itemize}
        \item Dietary TG and cholesterol packaged with apoB-48
        \item Secreted into lymphatics $\rightarrow$ thoracic duct $\rightarrow$ bloodstream
    \end{itemize}
    
    \item \textbf{Acquisition of apolipoproteins}
    \begin{itemize}
        \item HDL donates apoC-II and apoE to nascent chylomicrons
    \end{itemize}
    
    \item \textbf{Triglyceride hydrolysis}
    \begin{itemize}
        \item \textbf{Lipoprotein lipase (LPL)} on capillary endothelium
        \item Activated by apoC-II
        \item Releases fatty acids for tissue uptake
        \item Chylomicron shrinks to \textbf{chylomicron remnant}
    \end{itemize}
    
    \item \textbf{Remnant uptake by liver}
    \begin{itemize}
        \item ApoE recognized by hepatic remnant receptors
        \item Receptor-mediated endocytosis
    \end{itemize}
\end{enumerate}

\subsubsection{Endogenous Pathway (Hepatic Lipids)}

\begin{enumerate}
    \item \textbf{VLDL assembly} in liver
    \begin{itemize}
        \item Hepatic TG and cholesterol packaged with apoB-100
        \item Requires microsomal triglyceride transfer protein (MTP)
    \end{itemize}
    
    \item \textbf{VLDL secretion and modification}
    \begin{itemize}
        \item Acquires apoC-II and apoE from HDL
        \item LPL hydrolyzes triglycerides
        \item VLDL $\rightarrow$ IDL (intermediate-density lipoprotein)
    \end{itemize}
    
    \item \textbf{IDL fate}
    \begin{itemize}
        \item 50\% taken up by liver via apoE
        \item 50\% converted to LDL by hepatic lipase
    \end{itemize}
    
    \item \textbf{LDL metabolism}
    \begin{itemize}
        \item Main carrier of cholesterol to peripheral tissues
        \item Taken up via LDL receptor (recognizes apoB-100)
        \item Receptor-mediated endocytosis
    \end{itemize}
\end{enumerate}

\subsubsection{Reverse Cholesterol Transport (HDL Pathway)}

\begin{enumerate}
    \item \textbf{Nascent HDL formation}
    \begin{itemize}
        \item Liver and intestine secrete apoA-I
        \item Forms discoidal particles with phospholipids
    \end{itemize}
    
    \item \textbf{Cholesterol acquisition}
    \begin{itemize}
        \item \textbf{ABCA1 transporter} on peripheral cells effluxes cholesterol
        \item Free cholesterol transferred to nascent HDL
    \end{itemize}
    
    \item \textbf{Cholesterol esterification}
    \begin{itemize}
        \item \textbf{LCAT} (lecithin-cholesterol acyltransferase) on HDL
        \item Activated by apoA-I
        \item Esterifies cholesterol $\rightarrow$ moves to hydrophobic core
        \item Discoidal HDL $\rightarrow$ spherical HDL$_3$ $\rightarrow$ HDL$_2$
    \end{itemize}
    
    \item \textbf{Cholesterol delivery to liver}
    \begin{itemize}
        \item Direct: SR-B1 receptor on liver (selective uptake)
        \item Indirect: CETP transfers cholesteryl esters to VLDL/LDL
    \end{itemize}
\end{enumerate}

\begin{tcolorbox}[colback=blue!5, colframe=blue!70!black, title=\textbf{Clinical Significance of Lipoproteins}]
\textbf{Atherogenic:}
\begin{itemize}
    \item LDL ("bad cholesterol"): Deposits cholesterol in arterial walls
    \item Oxidized LDL: Taken up by macrophages $\rightarrow$ foam cells $\rightarrow$ atherosclerosis
\end{itemize}

\textbf{Atheroprotective:}
\begin{itemize}
    \item HDL ("good cholesterol"): Removes cholesterol from arteries
    \item Higher HDL associated with lower cardiovascular risk
\end{itemize}
\end{tcolorbox}

\subsection{Bile Acid Synthesis}

\begin{definition}
\textbf{Bile acids} are amphipathic molecules synthesized from cholesterol in the liver. They are essential for dietary lipid digestion and absorption, and represent the major route of cholesterol elimination.
\end{definition}

\textbf{Functions of Bile Acids:}
\begin{itemize}
    \item Emulsification of dietary fats
    \item Solubilization of cholesterol in bile
    \item Facilitation of fat-soluble vitamin absorption
    \item Elimination of excess cholesterol
\end{itemize}

\textbf{Primary Bile Acids (synthesized in liver):}
\begin{itemize}
    \item \textbf{Cholic acid} (3$\alpha$, 7$\alpha$, 12$\alpha$-trihydroxy)
    \item \textbf{Chenodeoxycholic acid} (3$\alpha$, 7$\alpha$-dihydroxy)
\end{itemize}

\textbf{Rate-Limiting Step:}
\[
\text{Cholesterol} \xrightarrow{\textbf{Cholesterol 7}\alpha\textbf{-hydroxylase (CYP7A1)}} \text{7}\alpha\text{-Hydroxycholesterol}
\]

\begin{itemize}
    \item Cytochrome P450 enzyme in liver ER
    \item Inhibited by bile acids (negative feedback)
    \item Induced by cholesterol
\end{itemize}

\textbf{Conjugation:}
\begin{itemize}
    \item Bile acids conjugated with glycine or taurine
    \item Forms glycocholic acid, taurocholic acid, etc.
    \item Conjugation increases water solubility and decreases pKa
    \item Conjugated bile acids = bile salts
\end{itemize}

\textbf{Secondary Bile Acids (formed by intestinal bacteria):}
\begin{itemize}
    \item Cholic acid $\rightarrow$ \textbf{Deoxycholic acid}
    \item Chenodeoxycholic acid $\rightarrow$ \textbf{Lithocholic acid}
    \item Formed by 7$\alpha$-dehydroxylation
\end{itemize}

\textbf{Enterohepatic Circulation:}
\begin{itemize}
    \item 95\% of bile acids reabsorbed in terminal ileum
    \item Return to liver via portal circulation
    \item Re-secreted into bile
    \item Pool cycles 6-10 times per day
    \item Only 5\% lost in feces (major route of cholesterol elimination)
\end{itemize}

\begin{tcolorbox}[colback=green!5, colframe=green!70!black, title=\textbf{Clinical Application: Bile Acid Sequestrants}]
\textbf{Drugs:} Cholestyramine, colestipol, colesevelam

\textbf{Mechanism:}
\begin{itemize}
    \item Bind bile acids in intestine
    \item Prevent reabsorption
    \item Increase fecal bile acid excretion
    \item Liver upregulates bile acid synthesis from cholesterol
    \item Upregulates LDL receptors to obtain more cholesterol
    \item Net effect: Decreased plasma LDL cholesterol
\end{itemize}
\end{tcolorbox}

\subsection{Steroid Hormone Synthesis}

Cholesterol serves as the precursor for all steroid hormones, synthesized in the adrenal cortex, gonads, and placenta.

\textbf{Rate-Limiting Step:}
\[
\text{Cholesterol} \xrightarrow{\textbf{Cholesterol desmolase (CYP11A1)}} \text{Pregnenolone}
\]
\begin{itemize}
    \item Located in inner mitochondrial membrane
    \item Also called side-chain cleavage enzyme (SCC)
    \item Removes 6-carbon side chain
    \item Stimulated by ACTH (adrenal) and LH (gonads)
    \item Requires StAR protein for cholesterol transport into mitochondria
\end{itemize}

\textbf{Classes of Steroid Hormones:}

\begin{center}
\begin{tabular}{lll}
\toprule
\textbf{Class} & \textbf{Examples} & \textbf{Site of Synthesis} \\
\midrule
Glucocorticoids & Cortisol & Adrenal cortex (zona fasciculata) \\
Mineralocorticoids & Aldosterone & Adrenal cortex (zona glomerulosa) \\
Androgens & Testosterone, DHEA & Adrenal cortex, testes \\
Estrogens & Estradiol & Ovaries, placenta \\
Progestogens & Progesterone & Ovaries, placenta \\
\bottomrule
\end{tabular}
\end{center}

\begin{tcolorbox}[colback=red!5, colframe=red!70!black, title=\textbf{Clinical Correlation: Congenital Adrenal Hyperplasia (CAH)}]
\textbf{Most common cause:} 21-Hydroxylase deficiency (CYP21A2)

\textbf{Pathophysiology:}
\begin{itemize}
    \item Cannot synthesize cortisol or aldosterone
    \item Precursors shunted to androgen pathway
    \item Elevated ACTH (no cortisol feedback)
    \item Adrenal hyperplasia
\end{itemize}

\textbf{Clinical features:}
\begin{itemize}
    \item Virilization of female infants (ambiguous genitalia)
    \item Salt-wasting crisis in severe cases
    \item Precocious puberty in males
    \item Elevated 17-hydroxyprogesterone (diagnostic)
\end{itemize}

\textbf{Treatment:} Glucocorticoid and mineralocorticoid replacement
\end{tcolorbox}

\section{Amino Acid Metabolism}

\subsection{Overview of Amino Acid Metabolism}

\begin{definition}
\textbf{Amino acid metabolism} encompasses the pathways for synthesis, degradation, and interconversion of amino acids, as well as the disposal of their nitrogen component.
\end{definition}

\textbf{Fates of Amino Acids:}
\begin{enumerate}
    \item \textbf{Protein synthesis}: Primary use
    \item \textbf{Precursors for specialized compounds}: Neurotransmitters, hormones, heme, nucleotides
    \item \textbf{Energy production}: Carbon skeletons enter TCA cycle
    \item \textbf{Gluconeogenesis}: Glucogenic amino acids
    \item \textbf{Ketogenesis}: Ketogenic amino acids
\end{enumerate}

\textbf{Sources of Amino Acids:}
\begin{itemize}
    \item Dietary protein digestion
    \item Intracellular protein turnover
    \item De novo synthesis (nonessential amino acids only)
\end{itemize}

\subsection{Essential vs. Nonessential Amino Acids}

\begin{definition}
\textbf{Essential amino acids} cannot be synthesized by the body and must be obtained from the diet. \textbf{Nonessential amino acids} can be synthesized endogenously.
\end{definition}

\begin{center}
\begin{tabular}{ll}
\toprule
\textbf{Essential (9)} & \textbf{Nonessential (11)} \\
\midrule
Histidine & Alanine \\
Isoleucine & Arginine* \\
Leucine & Asparagine \\
Lysine & Aspartate \\
Methionine & Cysteine* \\
Phenylalanine & Glutamate \\
Threonine & Glutamine \\
Tryptophan & Glycine \\
Valine & Proline \\
 & Serine \\
 & Tyrosine* \\
\bottomrule
\end{tabular}
\end{center}

*Conditionally essential: Required in diet under certain conditions (growth, illness)
\begin{tcolorbox}[colback=blue!5, colframe=blue!70!black, title=\textbf{Mnemonic for Essential Amino Acids}]
\textbf{"PVT TIM HALL"}
\begin{itemize}
    \item \textbf{P}henylalanine
    \item \textbf{V}aline
    \item \textbf{T}hreonine
    \item \textbf{T}ryptophan
    \item \textbf{I}soleucine
    \item \textbf{M}ethionine
    \item \textbf{H}istidine
    \item \textbf{A}rginine (conditionally essential)
    \item \textbf{L}eucine
    \item \textbf{L}ysine
\end{itemize}
\end{tcolorbox}

\subsection{Nitrogen Removal from Amino Acids}

Before amino acid carbon skeletons can be used for energy or biosynthesis, the amino group must be removed. This occurs through two main processes:

\subsubsection{Transamination}

\begin{definition}
\textbf{Transamination} is the transfer of an amino group from an amino acid to an $\alpha$-keto acid, producing a new amino acid and a new $\alpha$-keto acid.
\end{definition}

\textbf{General Reaction:}
\[
\text{Amino acid}_1 + \alpha\text{-Keto acid}_2 \xleftrightarrow{\text{Aminotransferase}} \alpha\text{-Keto acid}_1 + \text{Amino acid}_2
\]

\textbf{Key Features:}
\begin{itemize}
    \item \textbf{Reversible} reactions
    \item Require \textbf{pyridoxal phosphate (PLP)} — active form of vitamin B$_6$
    \item Do not result in net loss of nitrogen
    \item Collect amino groups onto glutamate
\end{itemize}

\textbf{Important Aminotransferases:}

\begin{enumerate}
    \item \textbf{Alanine Aminotransferase (ALT/GPT):}
    \[
    \text{Alanine} + \alpha\text{-Ketoglutarate} \xleftrightarrow{} \text{Pyruvate} + \text{Glutamate}
    \]
    \begin{itemize}
        \item Primarily in liver (cytosol)
        \item More specific for liver damage
        \item Important in glucose-alanine cycle
    \end{itemize}
    
    \item \textbf{Aspartate Aminotransferase (AST/GOT):}
    \[
    \text{Aspartate} + \alpha\text{-Ketoglutarate} \xleftrightarrow{} \text{Oxaloacetate} + \text{Glutamate}
    \]
    \begin{itemize}
        \item Found in liver, heart, muscle, kidney
        \item Present in both cytosol and mitochondria
        \item Elevated in liver disease and myocardial infarction
    \end{itemize}
\end{enumerate}

\begin{tcolorbox}[colback=red!5, colframe=red!70!black, title=\textbf{Clinical Correlation: Liver Function Tests}]
\textbf{ALT and AST} are measured as markers of liver damage:
\begin{itemize}
    \item \textbf{ALT > AST}: Suggests hepatocellular injury (hepatitis)
    \item \textbf{AST > ALT (ratio > 2:1)}: Suggests alcoholic liver disease
    \item \textbf{Both elevated}: General liver damage
\end{itemize}

\textbf{Note:} AST is also elevated in muscle damage and MI
\end{tcolorbox}

\subsubsection{Oxidative Deamination}

\begin{definition}
\textbf{Oxidative deamination} removes the amino group as free ammonia (NH$_3$), primarily from glutamate.
\end{definition}

\textbf{Key Enzyme: Glutamate Dehydrogenase (GDH)}
\[
\text{Glutamate} + \text{NAD}^+ (\text{or NADP}^+) + \text{H}_2\text{O} \xleftrightarrow{} \alpha\text{-Ketoglutarate} + \text{NH}_4^+ + \text{NADH}
\]

\textbf{Characteristics:}
\begin{itemize}
    \item Located in \textbf{mitochondrial matrix}
    \item Can use either NAD$^+$ or NADP$^+$
    \item \textbf{Reversible}: Direction depends on energy state
    \item Releases free ammonia for urea synthesis
\end{itemize}

\textbf{Regulation of Glutamate Dehydrogenase:}
\begin{itemize}
    \item \textbf{Activated by:} ADP, GDP (low energy signals)
    \item \textbf{Inhibited by:} ATP, GTP (high energy signals)
\end{itemize}

\subsection{Ammonia Transport}

Free ammonia is toxic, especially to the brain. Two main mechanisms safely transport nitrogen to the liver for urea synthesis:

\textbf{1. Glutamine Transport (Most Tissues):}
\[
\text{Glutamate} + \text{NH}_4^+ + \text{ATP} \xrightarrow{\text{Glutamine synthetase}} \text{Glutamine} + \text{ADP} + \text{P}_i
\]
\begin{itemize}
    \item Glutamine is non-toxic and can cross membranes
    \item In liver: Glutaminase releases NH$_4^+$ for urea synthesis
    \item In kidney: Glutaminase releases NH$_4^+$ for acid-base regulation
\end{itemize}

\textbf{2. Glucose-Alanine Cycle (Muscle):}

\begin{tcolorbox}[colback=green!5, colframe=green!70!black, title=\textbf{Glucose-Alanine Cycle}]
\textbf{In Muscle:}
\begin{enumerate}
    \item Amino acids undergo transamination with pyruvate
    \item Alanine is formed and released into blood
\end{enumerate}

\textbf{In Liver:}
\begin{enumerate}
    \item Alanine undergoes transamination to pyruvate
    \item Pyruvate used for gluconeogenesis
    \item Amino group enters urea cycle
    \item Glucose released to blood and returns to muscle
\end{enumerate}

\textbf{Purpose:} Transports nitrogen safely while recycling carbon
\end{tcolorbox}

\subsection{The Urea Cycle}

\begin{definition}
The \textbf{urea cycle} (Krebs-Henseleit cycle) converts toxic ammonia to non-toxic urea for excretion by the kidneys. It is the major pathway for nitrogen disposal in humans.
\end{definition}

\textbf{Location:} Liver (primarily periportal hepatocytes)
\begin{itemize}
    \item First two reactions: Mitochondrial matrix
    \item Remaining reactions: Cytosol
\end{itemize}

\textbf{Overall Reaction:}
\[
2\text{NH}_3 + \text{CO}_2 + 3\text{ATP} + \text{Aspartate} + 2\text{H}_2\text{O} \rightarrow \text{Urea} + 2\text{ADP} + 2\text{P}_i + \text{AMP} + \text{PP}_i + \text{Fumarate}
\]

\subsubsection{Reactions of the Urea Cycle}

\textbf{Reaction 1: Carbamoyl Phosphate Synthesis (Mitochondria)}
\[
\text{NH}_4^+ + \text{CO}_2 + 2\text{ATP} \xrightarrow{\textbf{CPS I}} \text{Carbamoyl phosphate} + 2\text{ADP} + \text{P}_i
\]

\begin{tcolorbox}[colback=yellow!10, colframe=yellow!70!black, title=\textbf{Carbamoyl Phosphate Synthetase I (CPS I)}]
\begin{itemize}
    \item \textbf{Rate-limiting enzyme} of urea cycle
    \item Located in mitochondrial matrix
    \item Requires \textbf{N-acetylglutamate (NAG)} as obligate activator
    \item Uses 2 ATP (one for activation, one for phosphorylation)
    \item Different from CPS II (cytosolic, pyrimidine synthesis)
\end{itemize}
\end{tcolorbox}

\textbf{Reaction 2: Citrulline Formation (Mitochondria)}
\[
\text{Carbamoyl phosphate} + \text{Ornithine} \xrightarrow{\text{Ornithine transcarbamoylase (OTC)}} \text{Citrulline} + \text{P}_i
\]
\begin{itemize}
    \item Citrulline is transported to cytosol via ornithine/citrulline antiporter
\end{itemize}

\textbf{Reaction 3: Argininosuccinate Formation (Cytosol)}
\[
\text{Citrulline} + \text{Aspartate} + \text{ATP} \xrightarrow{\text{Argininosuccinate synthetase}} \text{Argininosuccinate} + \text{AMP} + \text{PP}_i
\]
\begin{itemize}
    \item Second nitrogen enters from aspartate
    \item ATP $\rightarrow$ AMP (equivalent to 2 ATP)
\end{itemize}

\textbf{Reaction 4: Argininosuccinate Cleavage (Cytosol)}
\[
\text{Argininosuccinate} \xrightarrow{\text{Argininosuccinate lyase}} \text{Arginine} + \text{Fumarate}
\]
\begin{itemize}
    \item Fumarate can enter TCA cycle (links urea cycle to TCA cycle)
    \item This connection is called the "Krebs bicycle"
\end{itemize}

\textbf{Reaction 5: Urea Formation (Cytosol)}
\[
\text{Arginine} + \text{H}_2\text{O} \xrightarrow{\text{Arginase}} \text{Urea} + \text{Ornithine}
\]
\begin{itemize}
    \item Ornithine returns to mitochondria to continue cycle
    \item Urea diffuses into blood, excreted by kidneys
    \item Arginase is found almost exclusively in liver
\end{itemize}

\subsubsection{Regulation of the Urea Cycle}

\textbf{Short-term Regulation:}
\begin{itemize}
    \item \textbf{N-Acetylglutamate (NAG)}: Allosteric activator of CPS I
    \[
    \text{Glutamate} + \text{Acetyl-CoA} \xrightarrow{\text{NAG synthase}} \text{N-Acetylglutamate}
    \]
    \item NAG synthase is activated by arginine
    \item High protein diet $\rightarrow$ more amino acids $\rightarrow$ more arginine $\rightarrow$ more NAG $\rightarrow$ increased urea synthesis
\end{itemize}

\textbf{Long-term Regulation:}
\begin{itemize}
    \item High protein diet: Induces all urea cycle enzymes (10-20 fold)
    \item Starvation: Also induces enzymes (amino acids from protein breakdown)
    \item Low protein diet: Decreases enzyme levels
\end{itemize}

\subsubsection{Energy Cost of Urea Synthesis}

\begin{tcolorbox}[colback=blue!5, colframe=blue!70!black, title=\textbf{ATP Cost of Urea Synthesis}]
\textbf{Direct cost:}
\begin{itemize}
    \item CPS I: 2 ATP
    \item Argininosuccinate synthetase: 1 ATP $\rightarrow$ AMP + PP$_i$ (= 2 ATP equivalents)
\end{itemize}

\textbf{Total: 4 ATP equivalents per urea molecule}

\textbf{However:} Fumarate produced can generate $\sim$2.5 ATP via malate $\rightarrow$ OAA $\rightarrow$ NADH

\textbf{Net cost: $\sim$1.5 ATP per urea}
\end{tcolorbox}

\subsubsection{Clinical Correlations: Urea Cycle Disorders}

\begin{tcolorbox}[colback=red!5, colframe=red!70!black, title=\textbf{Urea Cycle Disorders}]
\textbf{General Features:}
\begin{itemize}
    \item All cause \textbf{hyperammonemia}
    \item Present in neonatal period or with protein load
    \item Symptoms: Vomiting, lethargy, seizures, coma
    \item Treatment: Low protein diet, nitrogen scavengers, dialysis
\end{itemize}

\textbf{Specific Disorders:}

\begin{center}
\begin{tabular}{lll}
\toprule
\textbf{Enzyme Deficiency} & \textbf{Accumulated Metabolite} & \textbf{Notes} \\
\midrule
CPS I deficiency & NH$_4^+$ & Decreased citrulline \\
OTC deficiency & Orotic acid, NH$_4^+$ & X-linked; most common \\
Argininosuccinate & Citrulline, NH$_4^+$ & \\
\quad synthetase & & \\
Argininosuccinate & Argininosuccinate & \\
\quad lyase & & \\
Arginase & Arginine & Least severe \\
NAG synthase & NH$_4^+$ & Treat with N-carbamylglutamate \\
\bottomrule
\end{tabular}
\end{center}
\end{tcolorbox}

\begin{tcolorbox}[colback=purple!5, colframe=purple!70!black, title=\textbf{Why Orotic Acid in OTC Deficiency?}]
In OTC deficiency:
\begin{enumerate}
    \item Carbamoyl phosphate accumulates in mitochondria
    \item Leaks into cytosol
    \item Enters pyrimidine synthesis pathway
    \item Excess orotic acid is produced and excreted
\end{enumerate}

\textbf{Distinguishing feature:} OTC deficiency has elevated orotic acid; CPS I deficiency does not
\end{tcolorbox}

\subsection{Amino Acid Carbon Skeleton Metabolism}

After nitrogen removal, amino acid carbon skeletons are converted to common metabolic intermediates.

\subsubsection{Classification by Metabolic Fate}

\begin{definition}
\textbf{Glucogenic amino acids} yield intermediates that can be converted to glucose via gluconeogenesis. \textbf{Ketogenic amino acids} yield acetyl-CoA or acetoacetate, which can form ketone bodies but not glucose.
\end{definition}

\begin{center}
\begin{tabular}{ll}
\toprule
\textbf{Category} & \textbf{Amino Acids} \\
\midrule
Purely Ketogenic & Leucine, Lysine \\
Both Glucogenic \& Ketogenic & Isoleucine, Phenylalanine, Threonine, \\
 & Tryptophan, Tyrosine \\
Purely Glucogenic & All others (13 amino acids) \\
\bottomrule
\end{tabular}
\end{center}

\begin{tcolorbox}[colback=blue!5, colframe=blue!70!black, title=\textbf{Mnemonic for Ketogenic Amino Acids}]
\textbf{"The Leu-Lys are purely ketogenic"}

For both ketogenic and glucogenic, remember: \textbf{"PITTT"}
\begin{itemize}
    \item \textbf{P}henylalanine
    \item \textbf{I}soleucine
    \item \textbf{T}hreonine
    \item \textbf{T}ryptophan
    \item \textbf{T}yrosine
\end{itemize}
\end{tcolorbox}

\subsubsection{Entry Points into Central Metabolism}

Amino acid carbon skeletons enter metabolism at seven points:

\begin{enumerate}
    \item \textbf{Pyruvate:} Alanine, Serine, Glycine, Cysteine, Threonine, Tryptophan
    
    \item \textbf{Acetyl-CoA:} Leucine, Isoleucine, Lysine, Threonine, Tryptophan
    
    \item \textbf{Acetoacetyl-CoA:} Leucine, Lysine, Phenylalanine, Tyrosine, Tryptophan
    
    \item \textbf{$\alpha$-Ketoglutarate:} Glutamate, Glutamine, Pr
oline, Histidine, Arginine
    
    \item \textbf{Succinyl-CoA:} Isoleucine, Valine, Methionine, Threonine
    
    \item \textbf{Fumarate:} Phenylalanine, Tyrosine, Aspartate
    
    \item \textbf{Oxaloacetate:} Aspartate, Asparagine
\end{enumerate}

\subsection{Metabolism of Specific Amino Acids}

\subsubsection{Branched-Chain Amino Acids (BCAAs)}

The branched-chain amino acids (leucine, isoleucine, valine) share common initial catabolic steps.

\textbf{Location:} Primarily in \textbf{skeletal muscle} (not liver)

\textbf{Step 1: Transamination}
\[
\text{BCAA} + \alpha\text{-ketoglutarate} \xrightarrow{\text{BCAA aminotransferase}} \alpha\text{-keto acid} + \text{Glutamate}
\]

\textbf{Step 2: Oxidative Decarboxylation (Rate-Limiting)}
\[
\alpha\text{-Keto acid} + \text{NAD}^+ + \text{CoA} \xrightarrow{\textbf{BCKDH}} \text{Acyl-CoA} + \text{CO}_2 + \text{NADH}
\]

\begin{tcolorbox}[colback=yellow!10, colframe=yellow!70!black, title=\textbf{Branched-Chain $\alpha$-Ketoacid Dehydrogenase (BCKDH)}]
\textbf{Structure:} Similar to pyruvate dehydrogenase and $\alpha$-ketoglutarate dehydrogenase

\textbf{Cofactors required:}
\begin{itemize}
    \item Thiamine pyrophosphate (TPP) — B$_1$
    \item Lipoic acid
    \item CoA (contains pantothenic acid — B$_5$)
    \item FAD — B$_2$
    \item NAD$^+$ — B$_3$
\end{itemize}

\textbf{Regulation:}
\begin{itemize}
    \item Activated by: BCAA excess, exercise
    \item Inhibited by: Phosphorylation (kinase), products (NADH, acyl-CoA)
\end{itemize}
\end{tcolorbox}

\begin{tcolorbox}[colback=red!5, colframe=red!70!black, title=\textbf{Clinical Correlation: Maple Syrup Urine Disease (MSUD)}]
\textbf{Defect:} Branched-chain $\alpha$-ketoacid dehydrogenase (BCKDH) deficiency

\textbf{Inheritance:} Autosomal recessive

\textbf{Biochemistry:}
\begin{itemize}
    \item Accumulation of BCAAs and their $\alpha$-keto acids
    \item Elevated leucine is most toxic (neurotoxic)
\end{itemize}

\textbf{Clinical features:}
\begin{itemize}
    \item Sweet, maple syrup odor of urine
    \item Poor feeding, vomiting
    \item Neurological deterioration, seizures
    \item Intellectual disability if untreated
    \item Onset: First week of life
\end{itemize}

\textbf{Treatment:}
\begin{itemize}
    \item Dietary restriction of BCAAs
    \item Thiamine supplementation (some variants are thiamine-responsive)
    \item Liver transplantation in severe cases
\end{itemize}
\end{tcolorbox}

\subsubsection{Phenylalanine and Tyrosine Metabolism}

\textbf{Phenylalanine Hydroxylation:}
\[
\text{Phenylalanine} + \text{O}_2 + \text{BH}_4 \xrightarrow{\textbf{Phenylalanine hydroxylase}} \text{Tyrosine} + \text{H}_2\text{O} + \text{BH}_2
\]

\begin{itemize}
    \item \textbf{BH$_4$} = tetrahydrobiopterin (cofactor)
    \item BH$_2$ is regenerated to BH$_4$ by dihydrobiopterin reductase
    \item This reaction makes tyrosine nonessential (if phenylalanine is adequate)
\end{itemize}

\begin{tcolorbox}[colback=red!5, colframe=red!70!black, title=\textbf{Clinical Correlation: Phenylketonuria (PKU)}]
\textbf{Defect:} Phenylalanine hydroxylase deficiency (most common) or BH$_4$ synthesis/regeneration defects

\textbf{Inheritance:} Autosomal recessive

\textbf{Biochemistry:}
\begin{itemize}
    \item Elevated phenylalanine in blood
    \item Decreased tyrosine (becomes essential)
    \item Alternative pathway produces phenylketones:
    \begin{itemize}
        \item Phenylpyruvate
        \item Phenyllactate
        \item Phenylacetate (musty odor)
    \end{itemize}
\end{itemize}

\textbf{Clinical features:}
\begin{itemize}
    \item Normal at birth
    \item Intellectual disability if untreated
    \item Fair skin, light hair, blue eyes (decreased melanin from low tyrosine)
    \item Musty/mousy body odor
    \item Eczema
    \item Seizures
\end{itemize}

\textbf{Screening:} Newborn screening (Guthrie test or tandem mass spectrometry)

\textbf{Treatment:}
\begin{itemize}
    \item Dietary phenylalanine restriction
    \item Tyrosine supplementation
    \item Avoid aspartame (contains phenylalanine)
    \item BH$_4$ supplementation for responsive variants
\end{itemize}

\textbf{Maternal PKU:} Pregnant women with PKU must maintain strict diet to prevent fetal damage (microcephaly, heart defects, intellectual disability)
\end{tcolorbox}

\textbf{Tyrosine Degradation Products:}
\begin{itemize}
    \item \textbf{Catecholamines}: Dopamine, norepinephrine, epinephrine
    \item \textbf{Thyroid hormones}: T$_3$, T$_4$
    \item \textbf{Melanin}: Skin and hair pigment
    \item \textbf{Fumarate and acetoacetate}: Energy metabolism
\end{itemize}

\begin{tcolorbox}[colback=red!5, colframe=red!70!black, title=\textbf{Clinical Correlation: Alkaptonuria}]
\textbf{Defect:} Homogentisate oxidase deficiency

\textbf{Inheritance:} Autosomal recessive

\textbf{Biochemistry:}
\begin{itemize}
    \item Accumulation of homogentisic acid (intermediate in Tyr degradation)
    \item Homogentisic acid oxidizes to dark pigment
\end{itemize}

\textbf{Clinical features:}
\begin{itemize}
    \item Dark urine (turns black on standing/exposure to air)
    \item Ochronosis (blue-black pigmentation of cartilage, connective tissue)
    \item Arthritis (especially spine, large joints)
    \item Dark earwax
    \item Generally benign; arthritis develops in adulthood
\end{itemize}

\textbf{Treatment:} Symptomatic; vitamin C may slow pigment deposition
\end{tcolorbox}

\subsubsection{Homocysteine Metabolism}

Homocysteine is a key intermediate in methionine metabolism with important clinical implications.

\textbf{Sources of Homocysteine:}
\[
\text{Methionine} \rightarrow \text{SAM} \rightarrow \text{SAH} \rightarrow \text{Homocysteine}
\]

\textbf{Fates of Homocysteine:}

\textbf{1. Remethylation to Methionine:}
\[
\text{Homocysteine} + \text{N}^5\text{-methyl-THF} \xrightarrow{\text{Methionine synthase}} \text{Methionine} + \text{THF}
\]
\begin{itemize}
    \item Requires \textbf{vitamin B$_{12}$} (methylcobalamin)
    \item Requires \textbf{folate} (as N$^5$-methyl-THF)
\end{itemize}

\textbf{2. Transsulfuration to Cysteine:}
\[
\text{Homocysteine} + \text{Serine} \xrightarrow{\text{Cystathionine } \beta\text{-synthase}} \text{Cystathionine}
\]
\[
\text{Cystathionine} \xrightarrow{\text{Cystathionase}} \text{Cysteine} + \alpha\text{-ketobutyrate}
\]
\begin{itemize}
    \item Both enzymes require \textbf{vitamin B$_6$} (PLP)
    \item Irreversible pathway
\end{itemize}

\begin{tcolorbox}[colback=red!5, colframe=red!70!black, title=\textbf{Clinical Correlation: Homocystinuria}]
\textbf{Causes of elevated homocysteine:}

\textbf{1. Cystathionine $\beta$-synthase (CBS) deficiency} (classic homocystinuria)
\begin{itemize}
    \item Most common cause
    \item Autosomal recessive
    \item Elevated: Homocysteine, methionine
    \item Decreased: Cysteine
\end{itemize}

\textbf{2. Vitamin B$_{12}$ or folate deficiency}
\begin{itemize}
    \item Impaired remethylation
    \item Elevated homocysteine, normal/low methionine
\end{itemize}

\textbf{3. MTHFR deficiency}
\begin{itemize}
    \item Cannot produce N$^5$-methyl-THF
    \item Impaired remethylation
\end{itemize}

\textbf{Clinical features of CBS deficiency:}
\begin{itemize}
    \item Marfanoid habitus (tall, long limbs)
    \item Downward lens dislocation (vs. upward in Marfan)
    \item Intellectual disability
    \item Thromboembolism (major cause of death)
    \item Osteoporosis
\end{itemize}

\textbf{Treatment:}
\begin{itemize}
    \item High-dose vitamin B$_6$ (some patients are responsive)
    \item Dietary methionine restriction
    \item Cysteine supplementation
    \item Folate and B$_{12}$ supplementation
    \item Betaine (alternative methyl donor)
\end{itemize}
\end{tcolorbox}

\subsection{Amino Acids as Precursors}

\subsubsection{Neurotransmitter Synthesis}

\begin{center}
\begin{tabular}{lll}
\toprule
\textbf{Amino Acid} & \textbf{Product} & \textbf{Key Enzyme/Cofactor} \\
\midrule
Tryptophan & Serotonin & Tryptophan hydroxylase, BH$_4$, B$_6$ \\
Tyrosine & Dopamine & Tyrosine hydroxylase, BH$_4$, B$_6$ \\
Tyrosine & Norepinephrine & Dopamine $\beta$-hydroxylase, vitamin C \\
Tyrosine & Epinephrine & PNMT, SAM \\
Glutamate & GABA & Glutamate decarboxylase, B$_6$ \\
Histidine & Histamine & Histidine decarboxylase, B$_6$ \\
\bottomrule
\end{tabular}
\end{center}

\subsubsection{Other Important Products}

\begin{itemize}
    \item \textbf{Glycine} $\rightarrow$ Heme, purines, glutathione, creatine
    \item \textbf{Glutamate} $\rightarrow$ Glutathione, GABA, glutamine
    \item \textbf{Arginine} $\rightarrow$ Nitric oxide (via NOS), creatine, urea
    \item \textbf{Tryptophan} $\rightarrow$ NAD$^+$/NADP$^+$ (niacin synthesis), melatonin
    \item \textbf{Methionine} $\rightarrow$ SAM (universal methyl donor), cysteine
    \item \textbf{Histidine} $\rightarrow$ Histamine
\end{itemize}

